\hypertarget{native-asyncawait-runtime-engine}{%
\chapter{NATIVE ASYNC/AWAIT RUNTIME
ENGINE}\label{native-asyncawait-runtime-engine}}

\hypertarget{pep-492-native-coroutines}{%
\subsection{11.1 PEP 492: Native
Coroutines}\label{pep-492-native-coroutines}}

Introduced in Python 3.5, \textbf{PEP 492} separated coroutines from
generators by introducing the \passthrough{\lstinline!async def!} and
\passthrough{\lstinline!await!} syntax. * \textbf{Native Coroutines
(\passthrough{\lstinline!PyCoroObject!})}: Functions declared with
\passthrough{\lstinline!async def!} return a native coroutine object
instead of a generator. Coroutines do not support standard iterator
methods (\passthrough{\lstinline!\_\_next\_\_!}); instead, they
implement the \passthrough{\lstinline!\_\_await\_\_!} method, which
returns an iterator used by the event loop to drive execution.

\hypertarget{the-event-loop-and-os-selectors}{%
\subsection{11.2 The Event Loop and OS
Selectors}\label{the-event-loop-and-os-selectors}}

The event loop is a single-threaded runtime manager that schedules and
executes concurrent operations. * \textbf{I/O Multiplexing}: The loop
utilizes the standard library \passthrough{\lstinline!selectors!} module
to monitor I/O sockets. This wraps high-performance OS-specific polling
mechanisms: - \passthrough{\lstinline!epoll!} (Linux) -
\passthrough{\lstinline!kqueue!} (macOS / BSD) -
\passthrough{\lstinline!select!} (Fallback platform wrapper) *
\textbf{Non-Blocking Scheduling}: When a task awaits an I/O operation
(like reading from a socket), the event loop registers the socket's file
descriptor with the OS selector and suspends the task. While the socket
is waiting, the loop executes other tasks. Once the selector signals
that the socket is ready, the event loop resumes the suspended task.

\hypertarget{task-and-future-state-machines}{%
\subsection{11.3 Task and Future State
Machines}\label{task-and-future-state-machines}}

\begin{itemize}
\tightlist
\item
  \textbf{Future}: Represents the eventual result of an asynchronous
  operation. It tracks state: \passthrough{\lstinline!PENDING!},
  \passthrough{\lstinline!CANCELLED!}, or
  \passthrough{\lstinline!FINISHED!}.
\item
  \textbf{Task}: A subclass of \passthrough{\lstinline!Future!} that
  wraps a coroutine object. The event loop drives the task by calling
  its \passthrough{\lstinline!step()!} method:

  \begin{enumerate}
  \def\labelenumi{\arabic{enumi}.}
  \tightlist
  \item
    \passthrough{\lstinline!step()!} invokes the coroutine's
    \passthrough{\lstinline!send(None)!} to run it.
  \item
    The coroutine executes until it awaits another future, which raises
    a \passthrough{\lstinline!YieldPoint!} and returns control to the
    loop.
  \item
    Once the awaited future completes, the loop calls
    \passthrough{\lstinline!step()!} again, passing the result back into
    the coroutine via \passthrough{\lstinline!send(result)!}.
  \end{enumerate}
\end{itemize}

\begin{lstlisting}
Task Lifecycle State Machine:
  [ PENDING ] ---> ( Scheduled on Loop ) ---> [ RUNNING (step called) ]
       |                                             |
       | <----( Yields on await / Pending Future )---+
       |
       +-------> ( Finished / Exception ) ------> [ FINISHED ]
\end{lstlisting}

\begin{center}\rule{0.5\linewidth}{0.5pt}\end{center}
